% ======================================================================
% 11-conclusion.tex (FULL REWRITE, Annals+++ / DIAMOND)
% PATH: src/conclusion/11-conclusion.tex
%
% Annals+++ contract:
%   - Conclusion = (i) Objects produced, (ii) Stability theorems (canonicality),
%                 (iii) Ledger as invariant, (iv) Interfaces (trace->zeta),
%                 (v) Seams (explicitly named, bounded, and localizable).
%   - Every claim is phrased as: invariant modulo explicit identity correction + ledger remainder.
%   - Ends with a STATUS_CERT style audit for portability.
% ======================================================================

\chapter{Conclusion and outlook}
\label{chap:11-conclusion}

\epigraph{A trace formula is useful only to the extent that its output is invariant under the admissible ways in which it can be implemented.}{}

% ======================================================================
\section{Summary of the main contributions}
\label{sec:11.summary}
% ======================================================================

The goal of this work was to build a localized trace framework in which the trace formula is treated as a \emph{stable
analytic object} rather than as a single formal expansion.  In microlocal implementations, auxiliary choices are
unavoidable---cutoff profiles, quantization conventions, partitions of unity, and (in noncompact settings) continuous-spectrum
regularizations.  The central thesis of the present work is that trace methods become \emph{portable} only after these
choices are systematically canonicalized and their residual influence is placed into a quantitative ledger.

The main outputs may be summarized as follows.

\begin{enumerate}
\item \textbf{Localized trace objects.}
A systematic construction of raw localized trace functionals $\mathcal L_\Gamma[h]$ associated with a fixed geometric input
$(X,\Gamma)$ and a fixed operator-theoretic setup (spectral resolution, regularization, microlocal projectors), together with
admissible test classes and rescaled families.

\item \textbf{Normalization layer (identity-sector control).}
A canonical subtraction map $\mathsf N_\chi(h)=h-h(0)\chi$ for an admissible $\chi\in C_c^\infty(\R)$ with $\chi\equiv1$ near $0$,
producing the normalized trace functional
\[
\mathcal T_\Gamma[h;\chi]:=\mathcal L_\Gamma[\mathsf N_\chi(h)].
\]

\item \textbf{Canonicality modulo explicit identity correction.}
A cutoff change formula and a normalization stability theorem asserting that all admissible modifications affect the normalized
trace only through an explicitly computable identity-sector correction plus a remainder controlled by a consolidated ledger.

\item \textbf{Consolidated error ledger as a structural invariant.}
A single quantitative object $\mathcal E_\Gamma[h;\chi]$ packaging all secondary sources of non-canonicity---quantization
remainders, refinement errors, truncation tails, and continuous-spectrum contributions---thereby providing a stable error budget
for every subsequent argument.

\item \textbf{Trace-to-zeta bridge (determinants, explicit formulae, positivity).}
A stable analytic interface transporting the normalized trace formalism to determinant-level encoders, explicit formulae, and
positivity criteria, with all noncanonical dependence confined to explicit identity corrections and ledger-bounded remainders.
\end{enumerate}

Taken together, these results yield a robust trace calculus in which spectral consequences are invariant statements
rather than artifacts of auxiliary microlocal choices.

% ======================================================================
\section{Normalization as canonicalization of trace objects}
\label{sec:11.normalization}
% ======================================================================

The raw trace functional $\mathcal L_\Gamma[h]$ contains an intrinsic identity-sector instability: admissible changes of the cutoff
profile modify the contribution carried by the constant mode $h(0)$, and this perturbation propagates unless it is removed at the
object level.

To isolate the stable content, fix $\chi\in C_c^\infty(\R)$ with $\chi\equiv1$ near $0$ and define
\[
\mathsf N_\chi(h):=h-h(0)\chi,
\qquad
\mathcal T_\Gamma[h;\chi]:=\mathcal L_\Gamma[\mathsf N_\chi(h)].
\]
The key structural property is that the dependence on $\chi$ is confined to a deterministic identity correction.

\begin{theorem}[Cutoff dependence is explicit]\label{thm:11.cutoff-explicit}
For any two admissible cutoffs $\chi_1,\chi_2$ and any admissible test function $h$,
\[
\mathcal T_\Gamma[h;\chi_1]-\mathcal T_\Gamma[h;\chi_2]
=
-h(0)\,\mathcal L_\Gamma[\chi_1-\chi_2].
\]
If the identity contribution in $\mathcal L_\Gamma$ is explicit, then the right-hand side is explicitly computable.
\end{theorem}

\noindent
Thus $\mathcal T_\Gamma[\cdot;\chi]$ may be treated as a canonical object modulo an explicit identity-sector correction.

% ======================================================================
\section{Stability under quantization conventions and microlocal refinement}
\label{sec:11.quantization}
% ======================================================================

A second obstruction to canonicity arises from the microlocal implementation: different admissible quantizations
(Weyl, Kohn--Nirenberg, etc.) and different microlocal partitions produce formally different kernels and trace expansions.
These differences must be shown to be lower order and quantitatively controlled.

To that end, the framework introduces discrepancy ledgers
\[
\mathcal E^{\mathrm{quant}}_\Gamma[h;\chi],
\qquad
\mathcal E^{\mathrm{ref}}_\Gamma[h;\chi],
\]
measuring, respectively, (i) changes induced by admissible quantization conventions and (ii) changes induced by admissible
refinements of microlocal covers.  The essential output is a uniform stability statement.

\begin{theorem}[Normalization stability under admissible auxiliary changes]\label{thm:11.norm-stability}
Let $(\chi,\Op,\Pi)$ and $(\chi',\Op',\Pi')$ be two admissible sets of auxiliary choices and let
$\mathcal T_\Gamma[h;\chi]$, $\mathcal T'_\Gamma[h;\chi']$ be the corresponding normalized traces.  Then for every admissible $h$,
\[
\mathcal T_\Gamma[h;\chi]-\mathcal T'_\Gamma[h;\chi']
=
\Delta_{\mathrm{id}}(h;\chi,\chi')
+
\Err_{\mathrm{norm}}(h),
\]
where $\Delta_{\mathrm{id}}(h;\chi,\chi')$ is an explicit identity-sector correction and
\[
|\Err_{\mathrm{norm}}(h)|
\le
\mathcal E_\Gamma[h;\chi]+\mathcal E'_\Gamma[h;\chi'].
\]
\end{theorem}

This converts a family of formally distinct microlocal trace realizations into a single stable invariant object equipped with an
explicit error budget.

% ======================================================================
\section{The consolidated ledger as a structural invariant}
\label{sec:11.ledger}
% ======================================================================

Rather than treating remainders informally, the present work makes the remainder calculus an explicit component of the object.
The consolidated ledger is defined by
\[
\mathcal E_\Gamma[h;\chi]
=
\mathcal E^{\mathrm{quant}}_\Gamma[h;\chi]
+
\mathcal E^{\mathrm{ref}}_\Gamma[h;\chi]
+
\mathcal E^{\mathrm{tail}}_\Gamma[h;\chi]
+
\mathcal E^{\mathrm{cont}}_\Gamma[h;\chi].
\]

This consolidation has two structural consequences.

\medskip
\noindent\textbf{(i) Invariance packaging.}
All noncanonical dependence of the localized trace method is isolated into a single controlled term, allowing later results to be
formulated solely in terms of the normalized trace $\mathcal T_\Gamma$ and its ledger $\mathcal E_\Gamma$.

\medskip
\noindent\textbf{(ii) Application interface.}
In semiclassical or large-parameter regimes, the usefulness of trace methods depends on a scale separation between main terms and
remainders.  The ledger is precisely the object that must be optimized (or shown negligible) to obtain meaningful quantitative
consequences.  In this sense, $\mathcal E_\Gamma$ is not an auxiliary estimate but an analytic invariant of the localization
architecture.

% ======================================================================
\section{From normalized traces to zeta-type objects}
\label{sec:11.zeta}
% ======================================================================

Chapter~\ref{chap:10-trace-to-zeta} recorded the standard analytic interfaces connecting trace formulae to zeta-type functions in a
form compatible with normalization and portability.

A key mechanism is the use of smoothed oscillatory test families
\[
h_{\tau,\Delta}(t)=W(t/\Delta)e^{-i\tau t},
\qquad W\in C_c^\infty(\R)\ \text{even},
\]
which produce the stable spectral probe
\[
\mathfrak D_\Gamma(\tau;\Delta):=\mathcal T_\Gamma[h_{\tau,\Delta}],
\]
admitting an explicit decomposition into geometric and spectral contributions with ledger-bounded error.

At the determinant level, trace identities reconstruct zeta-regularized determinants, and in settings admitting Euler products one
obtains logarithmic derivatives of zeta-type functions from geometric trace data.  Finally, positivity criteria arise canonically
from convolution squares: if $h^\ast(t)=\overline{h(-t)}$ and $h\star h^\ast$ is convolution, then
\[
\widehat{h\star h^\ast}(\xi)=|\widehat h(\xi)|^2\ge0,
\]
and the quadratic form
\[
\mathcal Q_\Gamma[h]:=\mathcal T_\Gamma[h\star h^\ast]
\]
yields sign constraints on spectral distributions.  This is the natural analytic route underlying Li--Weil type positivity
mechanisms.

The conceptual outcome is a stable bridge from microlocal trace analysis to zeta-type encoders in which all noncanonical dependence
is confined to explicit identity corrections and to the consolidated error ledger.

% ======================================================================
\section{Interpretation and methodological significance}
\label{sec:11.methodology}
% ======================================================================

Beyond its technical content, this work proposes a methodological viewpoint:

\begin{quote}
\emph{Trace formulae should be treated as canonical objects equipped with an error budget, not merely as formal expansions.}
\end{quote}

Classical presentations often hide the auxiliary choices required to implement microlocal localization.  While the resulting trace
identity may be correct, its later use becomes fragile unless invariance under those choices is explicitly tracked.
The normalization layer is therefore not a cosmetic adjustment: it is a structural upgrade that converts a trace formula into a
portable analytic instrument.

This is precisely what is required in problems where the gap between main terms and error terms is small and where stability under
admissible implementation changes is essential.

% ======================================================================
\section{Limitations and open analytic seams}
\label{sec:11.open}
% ======================================================================

The framework is designed to be portable across a range of trace-formula settings, but several analytic seams remain that must be
closed in concrete applications.  We record them explicitly.

\subsection{Sharp remainder bounds and ledger optimization}
\label{subsec:11.open.remainder}

The consolidated ledger packages errors canonically, but its effectiveness depends on sharp estimates.
In fine-scale regimes (e.g.\ positivity tests with small margins or extraction of refined spectral information) one needs sharper
bounds than those required for qualitative trace identities.  A primary direction for future work is scale-dependent optimization of
$\mathcal E_\Gamma[h;\chi]$, including explicit tradeoffs between cutoff width, microlocal resolution, and semiclassical scaling.

\subsection{Continuous spectrum and scattering determinants}
\label{subsec:11.open.cont}

In noncompact settings the continuous spectrum contributes scattering terms that must be incorporated into the normalization.
While the framework includes $\mathcal E^{\mathrm{cont}}_\Gamma[h;\chi]$, the explicit computation and minimization of this term
requires detailed control of scattering matrices and their analytic continuation, as well as stability under refinements and test
profile changes.

\subsection{Determinant-level completion}
\label{subsec:11.open.det}

Although the abstract determinant interface is recorded, a full implementation in a specific geometric model requires construction
of an operator family $K(s)$ whose determinant yields the relevant zeta object (Fredholm determinant realization).  This is classical
in Selberg-type settings but may be subtle beyond constant curvature.  A natural goal is a systematic dictionary between localized
trace functionals and Fredholm determinants in variable curvature and dynamical-zeta contexts.

% ======================================================================
\section{Final conclusion}
\label{sec:11.final}
% ======================================================================

The principal outcome of this work is a normalization framework that upgrades localized trace formulae into stable analytic objects.
The normalized trace functional $\mathcal T_\Gamma[h;\chi]$ is canonical modulo an explicit identity-sector correction, and its
dependence on admissible microlocal auxiliary choices is quantitatively controlled by a consolidated error ledger
$\mathcal E_\Gamma[h;\chi]$.

Consequently, subsequent spectral consequences---local Weyl laws, oscillatory trace observables, determinant identities, explicit
formulae, and positivity criteria---may be stated invariantly and applied without ambiguity, provided the relevant ledger bounds are
available in the regime of interest.

This provides a foundation for further investigations in spectral geometry and automorphic analysis in precisely those regimes where
stability and sharp error control are decisive.

% ======================================================================
\section*{Bibliographic notes}
% ======================================================================

The normalization viewpoint is consistent with classical microlocal trace theory and the Selberg trace formula tradition; see
\cite{Hejhal1983,Buser1992,Iwaniec2002}.  For semiclassical analysis and trace methods, see \cite{Zworski2012,DyatlovZworski2019}.
For spectral zeta functions and regularized determinants, see \cite{Seeley1967,RaySinger1971}.  For explicit formulae and zeta-related
positivity criteria, see \cite{Titchmarsh1986,Connes1999,Li1997}.

% ======================================================================
% Audit (Conclusion)
% ======================================================================

\section*{Audit (Conclusion)}\label{sec:11.audit}

\noindent\textbf{Objects produced.}
\begin{itemize}
\item Normalized trace $\mathcal T_\Gamma[h;\chi]$ (canonical modulo identity correction).
\item Consolidated ledger $\mathcal E_\Gamma[h;\chi]$ (quantization/refinement/tail/continuous-spectrum).
\item Stable bridges: oscillatory probes $\mathfrak D_\Gamma(\tau;\Delta)$; determinant encoder; positivity quadratic forms.
\end{itemize}

\noindent\textbf{Portability guarantees.}
\begin{itemize}
\item Identity-sector dependence explicit (cutoff change formula).
\item All other admissible auxiliary dependence absorbed into ledger-bounded remainder.
\end{itemize}

\noindent\textbf{Open seams (named).}
\begin{itemize}
\item Sharp ledger optimization in scale-dependent regimes.
\item Explicit continuous-spectrum/scattering control under refinement.
\item Determinant-level realization $K(s)$ beyond classical Selberg contexts.
\end{itemize}

\noindent\textbf{STATUS\_CERT.}
CLOSED (structural): conclusion stated entirely in terms of canonical objects + explicit identity correction + ledger.
NEXT(1): in the target geometric model, close a regime where $\mathcal E_\Gamma[h_{\tau,\Delta}]$ and $\mathcal E_\Gamma[h\star h^\ast]$
are negligible relative to the main terms needed for the intended positivity/determinant application.

% ======================================================================
% End of 11-conclusion.tex
% ======================================================================
